\subsection{P2P Marketplace}

\textbf{Peer-to-Peer (P2P) marketplace} là nền tảng số kết nối trực tiếp giữa những người sở hữu tài sản hoặc cung cấp dịch vụ với những người có nhu cầu sử dụng, mà không yêu cầu nền tảng đó phải sở hữu tài sản hay dịch vụ đó. 
Nó tạo môi trường để cá nhân tương tác, giao dịch, thuê mướn, mua bán hoặc chia sẻ sản phẩm/dịch vụ dựa trên niềm tin, đánh giá và quy tắc chung.

\subsubsection*{Các đặc điểm cốt lõi}

\begin{enumerate}[label=\arabic*.]
    \item \textbf{Nguồn cung do người dùng tạo ra:} Người sở hữu tài sản (ví dụ: nhà, xe, thiết bị) đăng tải thông tin lên nền tảng để cho thuê hoặc bán.

    \item \textbf{Giao dịch trực tiếp giữa các cá nhân:} Không qua trung gian sở hữu tài sản. 
    Nền tảng chỉ đóng vai trò cơ sở hạ tầng hỗ trợ (kết nối, thanh toán, đánh giá).

    \item \textbf{Cơ chế tin cậy (trust \& reputation):} 
    Bao gồm hệ thống đánh giá, phản hồi (ratings \& reviews), hồ sơ người dùng và xác thực danh tính nhằm đảm bảo cả người cho thuê và người thuê đều đáng tin cậy.

    \item \textbf{Tính an toàn và bảo mật:} 
    Hệ thống thanh toán bảo mật, cơ chế giải quyết tranh chấp, tuân thủ quy định pháp lý và bảo vệ quyền riêng tư của người dùng.
\end{enumerate}

\subsection{Các hình thức phổ biến của P2P Marketplace}

P2P marketplace có nhiều dạng, tùy theo loại tài sản hoặc dịch vụ được chia sẻ. 
Dưới đây là các hình thức phổ biến dựa theo tài liệu từ \textit{Flipkart Commerce Cloud}\cite{flipkartCCC2025}:

\begin{table}[h!]
\centering
\caption{Các hình thức phổ biến của P2P Marketplace}
\begin{tabular}{|p{3cm}|p{7cm}|p{3.5cm}|}
\hline
\textbf{Hình thức} & \textbf{Mô tả} & \textbf{Ví dụ} \\ \hline
\textbf{Sharing \newline Economy} & 
Cá nhân chia sẻ tài sản hoặc cung cấp dịch vụ cá nhân mà người khác cần sử dụng, theo thời gian hoặc theo dịch vụ. & 
Uber, Airbnb, các dịch vụ vệ sinh, dịch vụ nhỏ lẻ \\ \hline

\textbf{Crowdfunding} & 
Người khởi xướng dự án kêu gọi nhiều người góp vốn để thực hiện một sản phẩm hoặc một dự án chung. &
Kickstarter, Indiegogo \\ \hline

\textbf{Service \newline Marketplace} &
Người có kỹ năng cung cấp dịch vụ cho người cần: dạy học, coaching, thiết kế, sửa chữa, nghệ thuật ... &
Udemy, Fiverr, các nền tảng freelancer \\ \hline
\end{tabular}
\end{table}

\subsubsection*{Lợi ích chung của các mô hình P2P}

\begin{itemize}
    \item \textbf{Giảm chi phí khởi nghiệp:} Vì nền tảng không cần sở hữu lượng hàng tồn (inventory) lớn.
    \item \textbf{Hiệu ứng mạng (network effect):} Càng nhiều người dùng ở cả hai phía cung và cầu thì giá trị nền tảng càng tăng.
    \item \textbf{Khả năng mở rộng nhanh:} Nếu mô hình thị trường được thiết kế tốt và tính tin cậy được xây dựng vững chắc.
\end{itemize}

\subsubsection*{Thách thức chung}

\begin{itemize}
    \item Xây dựng niềm tin và kiểm soát chất lượng giữa các cá nhân.
    \item Đáp ứng quy định pháp lý, bảo hiểm, và trách nhiệm khi xảy ra sự cố.
    \item Đạt được số lượng người dùng đủ lớn để thị trường hoạt động hiệu quả (\textit{critical mass}).
    \item Thiết lập hệ thống thanh toán, xử lý tranh chấp và hỗ trợ người dùng.
\end{itemize}

\subsection{Sharing Economy }
Mô hình Sharing Economy (kinh tế chia sẻ) là một nhánh đặc thù trong hệ sinh thái P2P Marketplace, trong đó các cá nhân chia sẻ hoặc cho thuê lại tài sản, nguồn lực của mình thông qua một nền tảng trung gian số. Thay vì sở hữu riêng biệt, người dùng có thể tối ưu hóa giá trị sử dụng tài sản bằng cách chia sẻ quyền sử dụng tạm thời cho người khác.

Theo Flipkart Commerce Cloud \cite{flipkartCCC2025}, bản chất của kinh tế chia sẻ nằm ở việc chuyển từ quyền sở hữu (ownership) sang quyền sử dụng (access). Đây là bước chuyển lớn trong hành vi tiêu dùng hiện đại, nhấn mạnh việc sử dụng tài nguyên hiệu quả, giảm lãng phí và tạo ra lợi ích cho cả hai phía: người sở hữu và người sử dụng.

\subsubsection*{Đặc điểm nổi bật của mô hình Sharing Economy:}
\begin{itemize}
    \item Tối ưu hóa tài sản nhàn rỗi: Cho phép cá nhân khai thác tài sản chưa được sử dụng hết (ví dụ xe hơi, nhà ở, thiết bị).
    \item Tăng cường tính bền vững: Giảm nhu cầu sản xuất mới, từ đó giảm tiêu thụ năng lượng và phát thải.
    \item Cộng đồng người dùng dựa trên niềm tin: Các giao dịch dựa vào hệ thống đánh giá, xác thực và uy tín cá nhân.
    \item Hỗ trợ công nghệ nền: Sử dụng dữ liệu lớn (Big Data), định vị (GPS), thanh toán điện tử, và các API kết nối để đảm bảo trải nghiệm liền mạch.
\end{itemize}

\subsubsection*{Các lĩnh vực ứng dụng phổ biến:}
\begin{itemize}
    \item Vận chuyển: Uber, Grab, BlaBlaCar.
    \item Lưu trú: Airbnb, Luxstay.
    \item Tài chính cộng đồng: Lending Club, GoFundMe.
    \item Dịch vụ cá nhân: TaskRabbit, Fiverr.
\end{itemize}